Let $\Delta V$ be the voltage of the battery.

Before: A is in series with the pair B\,$\|$\,C. At the same voltage the pair lets more current through than one bulb (two paths), so in the voltage divider A takes the larger share: $\Delta V_A > \frac{1}{2}\,\Delta V$, and B and C each get $\Delta V - \Delta V_A < \frac{1}{2}\,\Delta V$.

After: no current flows through the open branch, so C goes off. A and B are now two identical bulbs in series, and they share the voltage equally: $\Delta V_A = \Delta V_B = \frac{1}{2}\,\Delta V$.

\begin{itemize}
\item A: its voltage falls from more than $\frac{1}{2}\,\Delta V$ to $\frac{1}{2}\,\Delta V$, its current falls too: A gets \result{\text{dimmer}}.
\item B: its voltage rises from less than $\frac{1}{2}\,\Delta V$ to $\frac{1}{2}\,\Delta V$: B gets \result{\text{brighter}}. A and B are now equally bright.
\item C: \result{\text{off}}.
\end{itemize}

The load now lets less current through, so the battery delivers less current. A typical mistake is to think that the current of C simply goes to B while A stays the same: the battery keeps a fixed voltage, not a fixed current.

A check with bulbs of constant resistance $R$: before, A has $\frac{2}{3}\,\Delta V$ and B has $\frac{1}{3}\,\Delta V$; after, each has $\frac{1}{2}\,\Delta V$. With $P = \Delta V^2/R$ the power of A falls from $\frac{4}{9}$ to $\frac{1}{4}$ of $\Delta V^2/R$, that of B rises from $\frac{1}{9}$ to $\frac{1}{4}$.
